% TODO make hw stylesheet
\documentclass[10pt]{scrartcl}
\usepackage[a4paper, total={6in, 8in}]{geometry}
\usepackage[usenames,dvipsnames,svgnames]{xcolor}
\usepackage[shortlabels]{enumitem}
\usepackage[framemethod=TikZ]{mdframed}
\usepackage{amsmath,amssymb,amsthm}
\usepackage[colorlinks]{hyperref}
\usepackage{multicol}
\usepackage{tabularx}

\hypersetup{
	colorlinks=true,
	linkcolor=blue,
	filecolor=magenta,      
	urlcolor=magenta,
	pdftitle={MAT 362 HW}
}

\usepackage{microtype}
\usepackage{mathtools}
\usepackage[headsepline]{scrlayer-scrpage}
\usepackage{thmtools}
\usepackage[textsize=footnotesize]{todonotes}

\mdfdefinestyle{mdbluebox}{roundcorner=10pt,innerbottommargin=9pt,
	linecolor=blue,backgroundcolor=TealBlue!5,}
\declaretheoremstyle[headfont=\sffamily\bfseries\color{MidnightBlue},
mdframed={style=mdbluebox},]{thmbluebox}

\mdfdefinestyle{mdbloobox}{frametitlefont=\bfseries,innerbottommargin=8pt,
	nobreak=true,backgroundcolor=TealBlue!5,linecolor=blue,}
\declaretheoremstyle[headfont=\bfseries\color{MidnightBlue},
mdframed={style=mdbloobox},headpunct={\\[3pt]},postheadspace=0pt,]{thmbloobox}

\mdfdefinestyle{mdredbox}{frametitlefont=\bfseries,innerbottommargin=8pt,
	nobreak=true,backgroundcolor=Salmon!5,linecolor=RawSienna,}
\declaretheoremstyle[headfont=\bfseries\color{RawSienna},
mdframed={style=mdredbox},headpunct={\\[3pt]},postheadspace=0pt,]{thmredbox}

\mdfdefinestyle{mdgreenbox}{linecolor=ForestGreen,backgroundcolor=ForestGreen!5,
	linewidth=2pt,rightline=false,leftline=true,topline=false,bottomline=false,}
\declaretheoremstyle[headfont=\bfseries\sffamily\color{ForestGreen!70!black},
mdframed={style=mdgreenbox},headpunct={ --- },]{thmgreenbox}

\mdfdefinestyle{mdorangebox}{linecolor=RedOrange,backgroundcolor=RedOrange!5,
	linewidth=2pt,rightline=false,leftline=true,topline=false,bottomline=false,}
\declaretheoremstyle[headfont=\bfseries\sffamily\color{RedOrange!70!black},
mdframed={style=mdorangebox},headpunct={ --- },]{thmorangebox}

\mdfdefinestyle{mdblackbox}{linecolor=black,backgroundcolor=RedViolet!5!gray!5,
	linewidth=3pt,rightline=false,leftline=true,topline=false,bottomline=false,}
\declaretheoremstyle[mdframed={style=mdblackbox}]{thmblackbox}

\declaretheoremstyle[spaceabove=6pt,spacebelow=6pt]{basehead}

\declaretheorem[style=basehead,name=Problem]{problem}
\declaretheorem[style=basehead,name=Problem,numbered=no]{problem*}
\declaretheorem[style=thmbloobox,name=Setup]{setup} % for config geo
\declaretheorem[style=thmredbox,name=Example,sibling=problem]{example}
\declaretheorem[style=thmredbox,name=$\bigstar$ Problem,sibling=problem]{niceprob}
\declaretheorem[style=thmbluebox,name=Theorem,sibling=problem]{theorem}
\declaretheorem[style=thmbluebox,name=Corollary,sibling=theorem]{corollary}
\declaretheorem[style=thmbluebox,name=Proposition,sibling=theorem]{prop}
\declaretheorem[style=thmbluebox,name=Lemma,numberwithin=problem]{lemma}
\declaretheorem[style=thmbluebox,name=Theorem,numbered=no]{theorem*}
\declaretheorem[style=thmbluebox,name=Lemma,numbered=no]{lemma*}
\declaretheorem[style=thmgreenbox,name=Claim,numberwithin=problem]{claim}
\declaretheorem[style=thmgreenbox,name=Claim,numbered=no]{claim*}
\declaretheorem[style=thmblackbox,name=Remark,numberwithin=problem]{remark}
\declaretheorem[style=thmblackbox,name=Remark,numbered=no]{remark*}
\declaretheorem[style=thmgreenbox,name=Definition,sibling=theorem]{definition}
\declaretheorem[style=thmgreenbox,name=Definition,numbered=no]{definition*}
\declaretheorem[style=thmorangebox,name=Warning,sibling=theorem]{warning}
\declaretheorem[style=thmorangebox,name=Warning,numbered=no]{warning*}
\declaretheorem[style=thmorangebox,name=Note,numbered=no]{note}
\declaretheorem[style=thmblackbox,name=Exercise,sibling=problem]{exercise}
\declaretheorem[style=thmblackbox,name=Exercise,numbered=no]{exercise*}
\declaretheorem[style=thmblackbox,name=Question,sibling=problem]{question}
\declaretheorem[style=thmblackbox,name=Question,numbered=no]{question*}

\addtolength{\textheight}{3.14cm}
\providecommand{\ol}{\overline}
\providecommand{\ul}{\underline}
\providecommand{\eps}{\varepsilon}
\providecommand{\half}{\frac{1}{2}}
\providecommand{\dang}{\measuredangle} %% Directed angle
\providecommand{\CC}{\mathbb C}
\providecommand{\FF}{\mathbb F}
\providecommand{\NN}{\mathbb N}
\providecommand{\QQ}{\mathbb Q}
\providecommand{\RR}{\mathbb R}
\providecommand{\ZZ}{\mathbb Z}
\providecommand{\dg}{^\circ}
\providecommand{\ii}{\item}
\providecommand{\setdiff}{\smallsetminus}
\providecommand{\alert}[1]{{\sffamily\textbf{\textcolor{blue}{#1}}}}
\providecommand{\inv}{^{-1}}
\providecommand{\mb}[1]{\mathbf{#1}}
\providecommand{\dif}{\;d}
\newcommand*{\horzbar}{\rule[.5ex]{2.5ex}{0.5pt}}

\reversemarginpar

\newenvironment{soln}{\begin{proof}[Solution]}{\end{proof}}
\newenvironment{walkthrough}{\noindent \textbf{Walkthrough.}}{\medskip}
\newlist{walk}{enumerate}{3}
\setlist[walk]{label=\bfseries (\alph*)}

\providecommand{\pow}{\operatorname{Pow}}
\providecommand{\vocab}[1]{{\textbf{\textcolor{ForestGreen}{#1}}}}
\providecommand{\lcm}{\operatorname{lcm}}
\providecommand{\abs}[1]{\left|#1\right|}

\usepackage{asymptote}
\begin{asydef}
	defaultpen(fontsize(10pt));
	unitsize(1cm);
	size(8cm); // set a reasonable default
	usepackage("amsmath");
	usepackage("amssymb");
	settings.tex="pdflatex";
	settings.outformat="pdf";
	// Replacement for olympiad+cse5 which is not standard
	import geometry;
	// recalibrate fill and filldraw for conics
	void filldraw(picture pic = currentpicture, conic g, pen fillpen=defaultpen, pen drawpen=defaultpen)
	{ filldraw(pic, (path) g, fillpen, drawpen); }
	void fill(picture pic = currentpicture, conic g, pen p=defaultpen)
	{ filldraw(pic, (path) g, p); }
	// some geometry
	pair foot(pair P, pair A, pair B) { return foot(triangle(A,B,P).VC); }
	pair orthocenter(pair A, pair B, pair C) { return orthocentercenter(A,B,C); }
	pair centroid(pair A, pair B, pair C) { return (A+B+C)/3; }
	// cse5 abbreviations
	path CP(pair P, pair A) { return circle(P, abs(A-P)); }
	path CR(pair P, real r) { return circle(P, r); }
	pair IP(path p, path q) { return intersectionpoints(p,q)[0]; }
	pair OP(path p, path q) { return intersectionpoints(p,q)[1]; }
	path Line(pair A, pair B, real a=0.6, real b=a) { return (a*(A-B)+A)--(b*(B-A)+B); }
	// cse5 more useful functions
	picture CC() {
		picture p=rotate(0)*currentpicture;
		currentpicture.erase();
		return p;
	}
	pair MP(Label s, pair A, pair B = plain.S, pen p = defaultpen) {
		Label L = s;
		L.s = "$"+s.s+"$";
		label(L, A, B, p);
		return A;
	}
	pair Drawing(Label s = "", pair A, pair B = plain.S, pen p = defaultpen) {
		dot(MP(s, A, B, p), p);
		return A;
	}
	path Drawing(path g, pen p = defaultpen, arrowbar ar = None) {
		draw(g, p, ar);
		return g;
	}
\end{asydef}

\usepackage{mathrsfs}

\ihead{\footnotesize MAT 362 HW06}
\ohead{\footnotesize Last compiled \today}

\title{\vspace{-2em}MAT 362 HW06}
\author{\large Aditya Pahuja}
\date{\large Due 03/13}
\begin{document}
\maketitle

\begin{problem*}3.2.2.
	Show that if a surface is tangent to a plane along a curve,
	then the points of this curve are either parabolic or planar.
\end{problem*}

\begin{soln}
	Suppose that the surface $S$ is tangent to the plane $P$
	along the curve $C$. As the point $p$ varies along $C$,
	$T_p(S) = P$ is constant, so $N(p)$ is also constant on $C$
	(since $N$ being differentiable means that it can't flip to $-N(p)$).
	Therefore, if for a given $p\in C$, the tangent vector to $C$ at $p$ is $w$,
	$dN_p(w)$ must be equal to zero, hence $dN_p$ isn't invertible
	and has determinant zero; thus the Gaussian curvature is zero
	at every point of $C$, which is exactly what it means
	for the points of the curve to be parabolic or planar.
\end{soln}

\begin{problem*}3.2.4.
	Assume a surface $S$ has the property that $|k_1|\le 1$, $|k_2|\le 1$
	everywhere. Is it true that the curvature $k$ of a curve on $S$
	also satisfies $|k|\le 1$?
\end{problem*}

\begin{soln}
	\framebox{No}, $|k|$ may be greater than $1$.
	Let $S$ be the $xy$-plane, which satisfies $|k_1|=|k_2|=0\le 1$ everywhere.
	Then, recall that the curvature of a circle with radius $r$
	is $\frac1r$ everywhere, so any circle in the $xy$-plane
	with radius less than $1$ will have curvature larger than $1$
	at each point on the circle. 
\end{soln}

\begin{problem*}3.2.5.
	Show that the mean curvature $H$ at $p\in S$ is given by
	\[H = \frac1\pi\int_0^\pi k_n(\theta)\dif\theta\]
	where $k_n(\theta)$ is the normal curvature at $p$
	along a direction making an angle $\theta$ with a fixed direction.
\end{problem*}

\begin{soln}
	Let $e_1$, $e_2$ be the principal directions at $p$.
	Letting $k_n(\theta)$ be the normal curvature at $p$
	along the direction making an angle of $\theta$ with $e_1$, we get
	\begin{align*}
		\int_0^\pi k_n(\theta)\dif\theta
		&= \int_0^\pi II_p(e_1\cos\theta + e_2\sin\theta)\dif\theta\\
		&= \int_0^\pi -\langle dN_p(e_1\cos\theta + e_2\sin\theta),
		e_1\cos\theta+e_2\sin\theta\rangle \dif\theta.
	\end{align*}
	Using the linearity of the inner product, we can write the integral as
	\[\int_0^\pi II_p(e_1)\cos^2\theta + II_p(e_2)\sin^2\theta
	- \sin\theta\cos\theta\cdot
	(\langle dN_p(e_1), e_2\rangle + \langle dN_p(e_2), e_1\rangle)\dif\theta.\]
	Then we can calculate
	\[\int_0^\pi \cos^2\theta\dif\theta =
	\int_0^\pi \frac{\cos(2\theta)+1}{2}\dif\theta
	= \frac\pi2 + \half\int_0^\pi\cos(2\theta)\dif\theta = \frac\pi2\]
	since
	\[\int_0^\pi\cos(2\theta)\dif\theta
	= \half\sin(2\pi) - \half\sin(2\cdot 0) = 0,\]
	and then
	\[\int_0^\pi \sin^2\theta\dif\theta
	= \int_0^\pi 1-\cos^2\theta\dif\theta
	= \pi - \int_0^\pi\cos^2\theta\dif\theta
	= \pi-\frac\pi2 = \frac\pi2.\]
	Additionally, by the identity $\sin(2\theta)=2\sin\theta\cos\theta$,
	\[\int_0^\pi\sin\theta\cos\theta\dif\theta
	= \half \int_0^\pi \sin(2\theta)\dif\theta
	= \half\cdot -\half(\cos(2\cdot\pi) - \cos(2\cdot 0)) = 0.\]
	Therefore,
	\[\frac1\pi\int_0^\pi II_p(e_1\cos\theta+e_2\sin\theta)\dif\theta
	= \frac1\pi\cdot\frac\pi2(II_p(e_1) + II_p(e_2))
	= \frac{k_1+k_2}{2}\]
	as desired.
\end{soln}

\begin{problem*}3.2.8.
	Describe the region of the unit sphere covered by the image
	of the Gauss map of the following surfaces:
	\begin{enumerate}[(a)]
		\ii Paraboloid of revolution $z=x^2+y^2$.
		\ii Hyperboloid of revolution $x^2+y^2-z^2=1$.
		\ii Catenoid $x^2+y^2=\cosh^2z$.
	\end{enumerate}
\end{problem*}

\begin{soln}
	Recall in general that if $f$ is a differentiable function
	for which $a$ is a regular value then the vector $(f_x,f_y,f_z)$
	is normal to the surface $f(x,y,z)=a$.
	
	(a) The paraboloid of revolution has an orientation given by
	\[N(x,y,x^2+y^2) = \frac{1}{\sqrt{4x^2+4y^2+1}}(2x,2y,-1).\]
	We will abbreviate this as $N(x,y)$ since the $z$-coordinate
	of a point on the surface is uniquely determined by $x$ and $y$.
	We show that the image of this map is the hemisphere $H$ of $S^2$
	with negative $z$-coordinate.
	Clearly the image is a subset of $H$, so we need to verify
	that every element of $H$ is in the image of $N$.
	That is, given $(a,b,c)\in S^2$ with $c<0$, there exist real $x$, $y$
	such that
	\[N(x,y) = \left(\frac{2x}{\sqrt{4x^2+4y^2+1}},
	\frac{2y}{\sqrt{4x^2+4y^2+1}},
	\frac{-1}{\sqrt{4x^2+4y^2+1}} \right)=(a,b,c).\]
	We can check that $x = \frac{-a}{2c}$,
	$y = \frac{-b}{2c}$ works by first noting that
	\[\sqrt{4x^2 + 4y^2 + 1} =
	\sqrt{4\left(\frac{-a}{2c}\right)^2 + 4\left(\frac{-b}{2c}\right)^2 + 1}
	= \sqrt{\frac{a^2+b^2+c^2}{c^2}} = \abs{\frac{1}{c}} = \frac{-1}{c},\]
	and then calculating
	\[N(x,y) = -c\left(2\cdot \frac{-a}{2c}, 2\cdot\frac{-b}{2c}, -1\right)
	= (a,b,c).\]
	
	(b) The hyperboloid of revolution has an orientation given by
	\[N(x,y,z) = \frac{1}{\sqrt{x^2+y^2+z^2}}(x,y,-z).\]
	We show that the image of $N$ is the subset of $S^2$
	strictly between the $45\dg$ and $-45\dg$ latitudes
	(i.e. the set of points $(a,b,c)\in S^2$ such that $|c|<\sqrt{a^2 + b^2}$).
	Clearly the image of $N$ is contained in this set
	because $x^2 + y^2 - z^2 = 1$ implies $z^2 < z^2 + 1\le x^2 + y^2$.
	Conversely, if $(a,b,c)\in S^2$ satisfies $|c|<\sqrt{a^2+b^2}$,
	then we wish to find a solution to $N(x,y,z)=(a,b,c)$.
	Take
	\[(x,y,z) = \left(\frac{a}{\sqrt{1-2c^2}}, \frac{b}{\sqrt{1-2c^2}},
	\frac{-c}{\sqrt{1-2c^2}}\right)\]
	(note that the denominators are real numbers because
	$c^2 < a^2 + b^2$ implies $c^2 < \half$).
	This satisfies
	\[x^2+y^2-z^2 = \frac{a^2+b^2-c^2}{1-2c^2} = \frac{1-c^2-c^2}{1-2c^2} = 1\]
	so the point lies on the hyperboloid, and
	$N(x,y,z)=(a,b,c)$ is obvious because $N(x,y,z)$ is just the unit vector
	in the direction of $(x,y,-z) = \frac{1}{\sqrt{1-2c^2}}(a,b,c)$.
	
	(c) The catenoid has an orientation given by
	\[N(x,y,z) = \frac{1}{\sqrt{x^2 + y^2 + (\cosh z\sinh z)^2}}
	(x,y,-\cosh z\sinh z).\]
	We show that the image of $N$ is all of $S^2$ except the poles $(0,0,\pm1)$.
	Clearly $N(x,y,z)$ cannot equal $(0,0,\pm1)$ because this requires
	$x=y=0$ but $\cosh^2z \ge 1$ for all $z$.
	Conversely, I claim that
	\[(x,y,z) = \left(\frac{a}{1-c^2}, \frac{b}{1-c^2},
	\operatorname{arctanh}(-c)\right)\]
	is a solution to $N(x,y,z)=(a,b,c)$ for all $(a,b,c)\in S^2$ with $|c|\ne1$.
	First we can compute
	\[-c = \tanh(z) = \frac{\sinh(z)}{\cosh(z)}
	= \frac{\sqrt{\cosh^2z - 1}}{\cosh z}\]
	which upon squaring and rearranging implies that $\cosh^2z=\frac{1}{1-c^2}$.
	This means that
	\[x^2 + y^2 = \frac{a^2 + b^2}{(1-c^2)^2} = \frac{1}{1-c^2}
	= \cosh^2z,\]
	so $(x,y,z)$ lies on the surface. Then,
	\begin{align*}
		N(x,y,z) &=
		\frac{1}{\sqrt{\cosh^2z + \cosh^2z\sinh^2z}}(x,y,-\cosh z\sinh z)\\
		&= \frac{1}{\cosh z\sqrt{1 + \sinh ^2 z}}(x,y,-\cosh z\sinh z)\\
		&= \frac{1}{\cosh^2 z}(x,y,-\cosh z\sinh z)\\
		&= \left(x(1-c^2), y(1-c^2), -\tanh z\right)\\
		&= (a,b,-(-c)) = (a,b,c)
	\end{align*}
	thus verifying that $N$ is a surjective map from the catenoid to $S^2
	\setdiff\{(0,0,1),(0,0,-1)\}$.
\end{soln}




























\end{document}